% =========================================================
% Axiomatic Equality
% =========================================================

\subsection{The Pure Language of Equality}

\begin{definitionbox}{Definition (Pure Language of Equality)}
\begin{definition}[Pure Language of Equality]\label{def:axiomatic-equality-language}
The \emph{pure language of equality}, denoted \(\mathcal L_{=}\), is the
language whose only distinguished logical relation is equality:
\[
  \mathcal L_{=}=\{=\}.
\]
It has no constant symbols, no function symbols, and no non-logical relation
symbols.
\end{definition}
\end{definitionbox}

\begin{remark*}[Standard quantified statement]
\[
  \mathsf{Const}_{\mathcal L_=}=\varnothing,\qquad
  \mathsf{Func}_{\mathcal L_=}=\varnothing,\qquad
  \mathsf{Rel}_{\mathcal L_=}^{\mathrm{nonlogical}}=\varnothing.
\]
\end{remark*}

\begin{remark*}[Interpretation]
The point of this language is to isolate the bare logical idea of identity
before adding subject-specific vocabulary such as membership, order, addition,
or distance.
\end{remark*}

\DefinitionalRoot

\begin{definitionbox}{Definition (Equality Structure)}
\begin{definition}[Equality Structure]\label{def:axiomatic-equality-structure}
An \(\mathcal L_{=}\)-structure is a pair
\[
  \mathcal M=(M,=^{\mathcal M})
\]
where \(M\neq\varnothing\) and equality is interpreted as the diagonal identity
relation
\[
  =^{\mathcal M}=\Delta_M
  =\{(x,x):x\in M\}\subseteq M\times M.
\]
\end{definition}
\end{definitionbox}

\begin{remark*}[Standard quantified statement]
\[
  \mathcal M=(M,=^{\mathcal M})
  \quad\text{with}\quad
  M\neq\varnothing
  \quad\text{and}\quad
  =^{\mathcal M}=\Delta_M.
\]
\end{remark*}

\begin{remark*}[Predicate reading]
\[
  \mathsf{EqualityStructure}(\mathcal M,M)
  \Longleftrightarrow
  M\neq\varnothing
  \;\wedge\;
  =^{\mathcal M}=\Delta_M.
\]
\end{remark*}

\begin{remark*}[Interpretation]
Once the domain \(M\) is chosen, the equality interpretation is fixed. Equality
is not an arbitrary binary relation selected independently of the domain.
\end{remark*}

\begin{dependencies}
\begin{itemize}
  \item \hyperref[def:axiomatic-equality-language]{Pure Language of Equality}
\end{itemize}
\end{dependencies}

\begin{definitionbox}{Definition (Distinctness)}
\begin{definition}[Distinctness]\label{def:distinctness}
The symbol \(\neq\) abbreviates non-identity:
\[
  x\neq y
  \quad\Longleftrightarrow\quad
  \neg(x=y).
\]
\end{definition}
\end{definitionbox}

\begin{remark*}[Standard quantified statement]
\[
  x\neq y
  \Longleftrightarrow
  \neg(x=y).
\]
\end{remark*}

\begin{remark*}[Interpretation]
Distinctness is not a new primitive relation. It is notation for the negation
of equality.
\end{remark*}

\begin{dependencies}
\begin{itemize}
  \item \hyperref[def:axiomatic-equality-language]{Pure Language of Equality}
\end{itemize}
\end{dependencies}

\subsection{Axioms}

\begin{axiombox}{\LraLogicalAxiomTitle{Axiom (Reflexivity of Equality)}}
\begin{axiom}[Reflexivity of Equality]\label{ax:axiomatic-equality-reflexivity}
Every object is identical to itself:
\[
  \forall x\,(x=x).
\]
\end{axiom}
\end{axiombox}

\begin{remark*}[Standard quantified statement]
\[
  \forall x\,(x=x).
\]
\end{remark*}

\begin{remark*}[Predicate reading]
\[
  \mathsf{Identical}(x,x).
\]
\end{remark*}

\begin{remark*}[Interpretation]
Reflexivity is the minimal identity law: every object may be replaced by
itself.
\end{remark*}

\begin{dependencies}
\begin{itemize}
  \item \hyperref[def:axiomatic-equality-language]{Pure Language of Equality}
\end{itemize}
\end{dependencies}

\begin{axiombox}{\LraLogicalAxiomTitle{Axiom Schema (Substitution of Identicals)}}
\begin{axiom}[Substitution of Identicals]\label{ax:axiomatic-equality-substitution}
For every formula \(\phi(x)\), if two objects are identical, then \(\phi\) has
the same truth value on both:
\[
  \forall x\,\forall y\,
  \bigl(x=y\to(\phi(x)\leftrightarrow\phi(y))\bigr).
\]
\end{axiom}
\end{axiombox}

\begin{remark*}[Standard quantified statement]
\[
  \forall x\,\forall y\,
  \bigl(x=y\to(\phi(x)\leftrightarrow\phi(y))\bigr).
\]
\end{remark*}

\begin{remark*}[Predicate reading]
\[
  x=y
  \Longrightarrow
  \bigl(\mathsf{PropertyHolds}(\phi,x)
  \leftrightarrow
  \mathsf{PropertyHolds}(\phi,y)\bigr).
\]
\end{remark*}

\begin{remark*}[Interpretation]
This is Leibniz's law: identical objects are indiscernible by any formula of
the language.

An application of this axiom begins by choosing a formula-context
\(\Phi(z)\), where \(z\) marks the place at which substitution is allowed.
If \(x=y\) and \(\Phi(x)\) has already been obtained, then substitution gives
\(\Phi(y)\). Other variables and other occurrences in the displayed expression
are held fixed as parameters of the chosen formula-context.

Thus substitution is not restricted to atomic predicates alone. It applies
inside relation statements, inside formulas involving function symbols, and
inside equations between terms, because each of these can be treated as a
formula \(\Phi(z)\). For example, to prove symmetry of equality one chooses
\(\Phi(z)\) to be \(z=x\). Reflexivity gives \(\Phi(x)\), namely \(x=x\);
from \(x=y\), substitution then gives \(\Phi(y)\), namely \(y=x\).
\end{remark*}

\begin{dependencies}
\begin{itemize}
  \item \hyperref[ax:axiomatic-equality-reflexivity]{Reflexivity of Equality}
\end{itemize}
\end{dependencies}

\subsection{Immediate Theorems}

\begin{theorembox}{Theorem (Symmetry of Equality)}
\begin{theorem}[Symmetry of Equality]\label{thm:axiomatic-equality-symmetry}
If \(x\) equals \(y\), then \(y\) equals \(x\):
\[
  \forall x\,\forall y\,(x=y\to y=x).
\]
\hyperref[prf:axiomatic-equality-symmetry]{\textit{Go to proof.}}
\end{theorem}
\end{theorembox}

\begin{remark*}[Standard quantified statement]
\[
  \forall x\,\forall y\,(x=y\to y=x).
\]
\end{remark*}

\begin{remark*}[Predicate reading]
\[
  \mathsf{Equal}(x,y)\Longrightarrow \mathsf{Equal}(y,x).
\]
\end{remark*}

\begin{remark*}[Interpretation]
Let \(\phi(z)\) be the formula \(z=x\). By reflexivity, \(\phi(x)\), namely
\(x=x\), is true. If \(x=y\), substitution transfers \(\phi\) from \(x\) to
\(y\), yielding \(y=x\).
\end{remark*}

\LeanFormalizes{thm:axiomatic-equality-symmetry}{lra-lean}{LRA.VolumeI.Identity.Axioms}{EqualitySymmetry}{checked}

\begin{dependencies}
\begin{itemize}
  \item \hyperref[ax:axiomatic-equality-reflexivity]{Reflexivity of Equality}
  \item \hyperref[ax:axiomatic-equality-substitution]{Substitution of Identicals}
\end{itemize}
\end{dependencies}

\begin{theorembox}{Theorem (Transitivity of Equality)}
\begin{theorem}[Transitivity of Equality]\label{thm:axiomatic-equality-transitivity}
If \(x=y\) and \(y=z\), then \(x=z\):
\[
  \forall x\,\forall y\,\forall z\,
  \bigl((x=y\land y=z)\to x=z\bigr).
\]
\hyperref[prf:axiomatic-equality-transitivity]{\textit{Go to proof.}}
\end{theorem}
\end{theorembox}

\begin{remark*}[Standard quantified statement]
\[
  \forall x\,\forall y\,\forall z\,
  \bigl((x=y\land y=z)\to x=z\bigr).
\]
\end{remark*}

\begin{remark*}[Predicate reading]
\[
  \mathsf{Equal}(x,y)\land\mathsf{Equal}(y,z)
  \Longrightarrow
  \mathsf{Equal}(x,z).
\]
\end{remark*}

\begin{remark*}[Interpretation]
Let \(\phi(w)\) be the formula \(x=w\). From \(x=y\), \(\phi(y)\) is true. From
\(y=z\), substitution transfers \(\phi\) from \(y\) to \(z\), giving \(x=z\).
\end{remark*}

\begin{dependencies}
\begin{itemize}
  \item \hyperref[ax:axiomatic-equality-substitution]{Substitution of Identicals}
\end{itemize}
\end{dependencies}

\begin{theorembox}{Theorem (Equality is an Equivalence Relation)}
\begin{theorem}[Equality is an Equivalence Relation]\label{thm:axiomatic-equality-equivalence-relation}
Equality is reflexive, symmetric, and transitive.
\[
  \mathsf{Reflexive}(=)\land\mathsf{Symmetric}(=)\land\mathsf{Transitive}(=).
\]
\hyperref[prf:axiomatic-equality-equivalence-relation]{\textit{Go to proof.}}
\end{theorem}
\end{theorembox}

\begin{remark*}[Standard quantified statement]
\[
  \bigl[\forall x(x=x)\bigr]
  \land
  \bigl[\forall x\,\forall y(x=y\to y=x)\bigr]
  \land
  \bigl[\forall x\,\forall y\,\forall z((x=y\land y=z)\to x=z)\bigr].
\]
\end{remark*}

\begin{remark*}[Predicate reading]
\[
  \mathsf{EquivalenceRelation}(=).
\]
\end{remark*}

\begin{remark*}[Interpretation]
Equality is the first and strictest equivalence relation: it identifies only an
object with itself.
\end{remark*}

\begin{dependencies}
\begin{itemize}
  \item \hyperref[ax:axiomatic-equality-reflexivity]{Reflexivity of Equality}
  \item \hyperref[thm:axiomatic-equality-symmetry]{Symmetry of Equality}
  \item \hyperref[thm:axiomatic-equality-transitivity]{Transitivity of Equality}
\end{itemize}
\end{dependencies}

\subsection{Congruence and Well-Defined Operations}

\begin{theorembox}{Theorem (Function Congruence)}
\begin{theorem}[Function Congruence]\label{thm:axiomatic-equality-function-congruence}
For any unary function symbol \(f\),
\[
  \forall x\,\forall y\,(x=y\to f(x)=f(y)).
\]
More generally, equal inputs in each coordinate yield equal outputs.
\hyperref[prf:axiomatic-equality-function-congruence]{\textit{Go to proof.}}
\end{theorem}
\end{theorembox}

\begin{remark*}[Standard quantified statement]
\[
  \forall x\,\forall y\,(x=y\to f(x)=f(y)).
\]
\end{remark*}

\begin{remark*}[Predicate reading]
\[
  \mathsf{Equal}(x,y)\Longrightarrow \mathsf{Equal}(f(x),f(y)).
\]
\end{remark*}

\begin{remark*}[Interpretation]
Function congruence is the logical reason operations are well-defined: applying
the same operation to identical inputs cannot produce different outputs.
\end{remark*}

\begin{dependencies}
\begin{itemize}
  \item \hyperref[ax:axiomatic-equality-substitution]{Substitution of Identicals}
\end{itemize}
\end{dependencies}

\begin{theorembox}{Theorem (Relation Congruence)}
\begin{theorem}[Relation Congruence]\label{thm:axiomatic-equality-relation-congruence}
For any relation symbol \(R\),
\[
  \forall x\,\forall y\,
  \bigl(x=y\to(R(x,z)\leftrightarrow R(y,z))\bigr),
\]
with the analogous rule in every coordinate.
\hyperref[prf:axiomatic-equality-relation-congruence]{\textit{Go to proof.}}
\end{theorem}
\end{theorembox}

\begin{remark*}[Standard quantified statement]
\[
  \forall x\,\forall y\,
  \bigl(x=y\to(R(x,z)\leftrightarrow R(y,z))\bigr).
\]
\end{remark*}

\begin{remark*}[Predicate reading]
\[
  \mathsf{Equal}(x,y)\Longrightarrow
  \bigl(\mathsf{Related}(R,x,z)\leftrightarrow\mathsf{Related}(R,y,z)\bigr).
\]
\end{remark*}

\begin{remark*}[Interpretation]
Relation congruence says identical objects have the same relational behavior
with respect to every relation expressible in the language.
\end{remark*}

\begin{dependencies}
\begin{itemize}
  \item \hyperref[ax:axiomatic-equality-substitution]{Substitution of Identicals}
\end{itemize}
\end{dependencies}

\subsection{Definitions Unlocked by Equality}

\begin{definitionbox}{Definition (Unique Existence)}
\begin{definition}[Unique Existence]\label{def:axiomatic-equality-unique-existence}
For a formula \(\psi(x)\), the expression \(\exists!x\,\psi(x)\) means
\[
  \exists!x\,\psi(x)
  \quad\Longleftrightarrow\quad
  \exists x\bigl(\psi(x)\land
  \forall y(\psi(y)\to y=x)\bigr).
\]
\end{definition}
\end{definitionbox}

\begin{remark*}[Standard quantified statement]
\[
  \exists!x\,\psi(x)
  \Longleftrightarrow
  \exists x\bigl(\psi(x)\land
  \forall y(\psi(y)\to y=x)\bigr).
\]
\end{remark*}

\begin{remark*}[Predicate reading]
\[
  \mathsf{UniqueWitness}(\psi)
  \Longleftrightarrow
  \text{there is exactly one object satisfying }\psi.
\]
\end{remark*}

\begin{remark*}[Interpretation]
Uniqueness is existence plus an equality test: every other witness must be the
same object as the chosen witness.
\end{remark*}

\begin{dependencies}
\begin{itemize}
  \item \hyperref[def:distinctness]{Distinctness}
  \item \hyperref[thm:axiomatic-equality-equivalence-relation]{Equality is an Equivalence Relation}
\end{itemize}
\end{dependencies}

\begin{definitionbox}{Definition (Two-Element Bounds)}
\begin{definition}[Two-Element Bounds]\label{def:two-element-bounds}
Equality allows finite-size constraints to be expressed internally:
\[
  \text{at least two:}\quad
  \exists x\,\exists y\,\neg(x=y),
\]
and
\[
  \text{at most two:}\quad
  \forall x\,\forall y\,\forall z\,
  (x=y\lor y=z\lor x=z).
\]
\end{definition}
\end{definitionbox}

\begin{remark*}[Standard quantified statement]
\[
  \exists x\,\exists y\,\neg(x=y)
  \qquad\text{and}\qquad
  \forall x\,\forall y\,\forall z\,
  (x=y\lor y=z\lor x=z).
\]
\end{remark*}

\begin{remark*}[Predicate reading]
\[
  \mathsf{AtLeastTwo}(M)
  \qquad\text{and}\qquad
  \mathsf{AtMostTwo}(M).
\]
\end{remark*}

\begin{remark*}[Interpretation]
Before number theory is available, equality already lets the language compare
objects and express small cardinal constraints.
\end{remark*}

\begin{dependencies}
\begin{itemize}
  \item \hyperref[def:distinctness]{Distinctness}
\end{itemize}
\end{dependencies}

\subsection{Bridge to First-Order Logic with Equality}

\begin{remark*}[Exposition]
Plain first-order logic need not include equality. Hereafter, Predicate Logic
is read as first-order logic with equality: the ordinary first-order syntax
plus the identity symbol governed by the axioms above.
\end{remark*}
